\documentclass[12pt]{article}
\usepackage{amsmath,amssymb,amsfonts}
\usepackage[margin=1in]{geometry}

\begin{document}

\textbf{Chapter 6, Problem 7.} Analyze the function $z^{1/n}$ in terms of the Riemann sheet concepts discussed in Section 6.2.

\bigskip
\textbf{Solution.}

Write $z = re^{i\theta}$ where $r > 0$ and $\theta = \arg z$. Then:
\[
z^{1/n} = r^{1/n} e^{i\theta/n}.
\]

However, $\theta$ is only determined up to multiples of $2\pi$. If we write $\theta \to \theta + 2k\pi$ for integer $k$, we get:
\[
z^{1/n} = r^{1/n} e^{i(\theta + 2k\pi)/n}, \qquad k = 0, 1, 2, \ldots, n-1.
\]

This gives $n$ distinct values for each $z \neq 0$, so $z^{1/n}$ is an $n$-valued function.

\textbf{Branch points and branch cuts.} The point $z = 0$ is a branch point: if we continuously vary $\theta$ from $0$ to $2\pi$ (going once around the origin), the function value changes from $r^{1/n}e^{i\theta/n}$ to $r^{1/n}e^{i(\theta+2\pi)/n}$, which is a different branch. After going around $n$ times ($\theta \to \theta + 2n\pi$), we return to the original value. Hence $z = 0$ is a branch point of order $n$.

The point at infinity is also a branch point, since $w = 1/z$ and $w^{1/n}$ has a branch point at $w = 0$.

\textbf{Riemann surface construction.} To make $z^{1/n}$ single-valued, we construct an $n$-sheeted Riemann surface:

\begin{enumerate}
\item Take $n$ copies (sheets) of the complex plane, each cut along the positive real axis (or any ray from the origin to infinity).

\item Label the sheets $k = 0, 1, \ldots, n-1$. On sheet $k$, define
\[
f_k(z) = r^{1/n} e^{i(\theta + 2k\pi)/n}, \qquad 0 \leq \theta < 2\pi.
\]

\item Join the sheets: the upper edge of the cut on sheet $k$ is joined to the lower edge on sheet $k+1$ (mod $n$). That is, crossing the cut from sheet $k$ moves to sheet $k+1$, and crossing from sheet $n-1$ returns to sheet $0$.
\end{enumerate}

On this $n$-sheeted Riemann surface, $z^{1/n}$ is a single-valued, analytic function. Each sheet carries one of the $n$ branches, and analytic continuation around the origin cycles through all $n$ sheets.

For example, for $n = 2$ (the square root), we get a two-sheeted surface with branch points at $0$ and $\infty$, and the function $\sqrt{z}$ is single-valued and analytic on this surface. For $n = 3$, three sheets are needed, and so on.

The function $z^{1/n}$ is analytic everywhere on each sheet except at the branch point $z = 0$ and along the branch cut. On the full Riemann surface, it is analytic everywhere except at $z = 0$ and $z = \infty$.

\end{document}
